% GENERATED FILE. Edit canonical lessons, then run npm run generate:books.
% canonical-source-hash: fnv1a64:f13b115d5b5db908
% canonical-lessons: MW-C230-balno-scorch, MW-C230-dhalno, MW-C230-dharti, MW-C230-ak, MW-C230-kato

\chapter{Burning Sand, Setting Sun}
\label{ch:mw-burning-sand-setting-sun}
\hlchaptermodality{230}

\begin{chapteropening}
\textbf{By the end of this chapter:} \emph{I can describe the scorching ground and the sun going down.}

Complete the last lesson of chapter 230: I can describe the scorching ground and the sun going down.
\end{chapteropening}

\section[\texorpdfstring{\mw{बळणो} (baḷṇo)}{बळणो (baḷṇo)}]{\mw{बळणो} (baḷṇo) — to be scorching; to burn (by itself)}

\label{lesson:MW-C230-balno-scorch}



\begin{quote}
Before the new one: say the Marwadi for white, then the Marwadi for to shine.
\end{quote}

\subsection*{You'll want to know: \mw{बळणो}}
\textbf{\mw{बळणो}} — \emph{baḷṇo} — \textquotedblleft{}to be scorching; to burn (by itself)\textquotedblright{}.

\begin{quote}
\mw{दोपार} \mw{मांय} \mw{रेत} \mw{बळै} \mw{है।}
\end{quote}

\begin{quote}
\emph{dopār mā̃y ret baḷai hai.}
\end{quote}

\begin{quote}
At midday the sand is scorching.
\end{quote}

Past (it …): \mw{वो} \mw{दीयो} \mw{बळ्यो}, \emph{vo dīyo baḷyo}.

\begin{cousinweb}
\textbf{\mw{बळणो}} (\emph{baḷṇo}) is inherited from a reconstructed \emph{*dvalati}, \textquotedblleft{}burns\textquotedblright{}; CDIAL lists the Old Marwari verb \emph{balaï}, \textquotedblleft{}burns\textquotedblright{} (CDIAL 6654). \textbf{\mw{बाळणो}}, to burn something, is its partner.
\end{cousinweb}

\begin{grammarlens}[title={What you've built}]
One more word for this chapter.
\end{grammarlens}

\subsection*{Your turn}
\begin{itemize}
\raggedright
  \item \emph{Say it:} \emph{baḷṇo}
  \item \emph{Say it:} \emph{baḷṇo}, once more
  \item \emph{Say it:} \emph{dopār mā̃y ret baḷai hai.} — At midday the sand is scorching.
  \item \emph{Recall:} say the Marwadi for white, then the Marwadi for to shine, then say \emph{baḷṇo} again
\end{itemize}

\begin{tcolorbox}[breakable,colback=teal!4,colframe=teal!35!black,title={Before you move on}]
What does \mw{बळणो} mean? (\textquotedblleft{}To be scorching; to burn (by itself)\textquotedblright{}.) Say it once more.
\end{tcolorbox}
\section[\texorpdfstring{\mw{ढळणो} (ḍhaḷṇo)}{ढळणो (ḍhaḷṇo)}]{\mw{ढळणो} (ḍhaḷṇo) — to go down (of the sun); to slope}

\label{lesson:MW-C230-dhalno}



\begin{quote}
Before the new one: say the Marwadi for to shine, then the Marwadi for to be scorching.
\end{quote}

\subsection*{You'll want to know: \mw{ढळणो}}
\textbf{\mw{ढळणो}} — \emph{ḍhaḷṇo} — \textquotedblleft{}to go down (of the sun); to slope\textquotedblright{}.

\begin{quote}
\mw{सूरज} \mw{आथण} \mw{ढळै} \mw{है।}
\end{quote}

\begin{quote}
\emph{sūraj āthaṇ ḍhaḷai hai.}
\end{quote}

\begin{quote}
The sun goes down in the evening.
\end{quote}

Past (it …): \mw{वो} \mw{सूरज} \mw{ढळ्यो}, \emph{vo sūraj ḍhaḷyo}.

\begin{cousinweb}
\textbf{\mw{ढळणो}} (\emph{ḍhaḷṇo}) is inherited from a reconstructed \emph{*ḍhalati}, \textquotedblleft{}bends over, falls\textquotedblright{}; CDIAL lists the Old Marwari verb \emph{ḍhalaï}, \textquotedblleft{}falls, declines\textquotedblright{} (CDIAL 5581).
\end{cousinweb}

\begin{grammarlens}[title={What you've built}]
One more word for this chapter.
\end{grammarlens}

\subsection*{Your turn}
\begin{itemize}
\raggedright
  \item \emph{Say it:} \emph{ḍhaḷṇo}
  \item \emph{Say it:} \emph{ḍhaḷṇo}, once more
  \item \emph{Say it:} \emph{sūraj āthaṇ ḍhaḷai hai.} — The sun goes down in the evening.
  \item \emph{Recall:} say the Marwadi for to shine, then the Marwadi for to be scorching, then say \emph{ḍhaḷṇo} again
\end{itemize}

\begin{tcolorbox}[breakable,colback=teal!4,colframe=teal!35!black,title={Before you move on}]
What does \mw{ढळणो} mean? (\textquotedblleft{}To go down (of the sun); to slope\textquotedblright{}.) Say it once more.
\end{tcolorbox}
\section[\texorpdfstring{\mw{धरती} (dhartī)}{धरती (dhartī)}]{\mw{धरती} (dhartī) — the earth; the ground}

\label{lesson:MW-C230-dharti}



\begin{quote}
Before the new one: say the Marwadi for to be scorching, then the Marwadi for to go down (of the sun).
\end{quote}

\subsection*{You'll want to know: \mw{धरती}}
\textbf{\mw{धरती}} — \emph{dhartī} — \textquotedblleft{}the earth; the ground\textquotedblright{}.

\begin{quote}
\mw{आज} \mw{धरती} \mw{ताती} \mw{है।}
\end{quote}

\begin{quote}
\emph{āj dhartī tātī hai.}
\end{quote}

\begin{quote}
Today the ground is hot.
\end{quote}

A feminine noun.

\begin{cousinweb}
\textbf{\mw{धरती}} (\emph{dhartī}) is inherited from Sanskrit \emph{dharitrī}, \textquotedblleft{}she who carries\textquotedblright{}, a name for the earth; CDIAL lists Old Marwari \emph{dharatī}, \textquotedblleft{}the earth\textquotedblright{} (CDIAL 6750).
\end{cousinweb}

\begin{grammarlens}[title={What you've built}]
One more word for this chapter.
\end{grammarlens}

\subsection*{Your turn}
\begin{itemize}
\raggedright
  \item \emph{Say it:} \emph{dhartī}
  \item \emph{Say it:} \emph{dhartī}, once more
  \item \emph{Say it:} \emph{āj dhartī tātī hai.} — Today the ground is hot.
  \item \emph{Recall:} say the Marwadi for to be scorching, then the Marwadi for to go down (of the sun), then say \emph{dhartī} again
\end{itemize}

\begin{tcolorbox}[breakable,colback=teal!4,colframe=teal!35!black,title={Before you move on}]
What does \mw{धरती} mean? (\textquotedblleft{}The earth; the ground\textquotedblright{}.) Say it once more.
\end{tcolorbox}
\section[\texorpdfstring{\mw{आक} (āk)}{आक (āk)}]{\mw{आक} (āk) — swallow-wort (a desert shrub)}

\label{lesson:MW-C230-ak}



\begin{quote}
Before the new one: say the Marwadi for to go down (of the sun), then the Marwadi for the earth.
\end{quote}

\subsection*{You'll want to know: \mw{आक}}
\textbf{\mw{आक}} — \emph{āk} — \textquotedblleft{}swallow-wort (a desert shrub)\textquotedblright{}.

\begin{quote}
\mw{रेत} \mw{मांय} \mw{आक} \mw{है।}
\end{quote}

\begin{quote}
\emph{ret mā̃y āk hai.}
\end{quote}

\begin{quote}
There's swallow-wort in the sand.
\end{quote}

A masculine noun.

\begin{cousinweb}
\textbf{\mw{आक}} (\emph{āk}) is inherited from Sanskrit \emph{arka}, the name of the plant \emph{Calotropis}; CDIAL lists Marwari \emph{āk} (CDIAL 625). It grows wild in the sand.
\end{cousinweb}

\begin{grammarlens}[title={What you've built}]
One more word for this chapter.
\end{grammarlens}

\subsection*{Your turn}
\begin{itemize}
\raggedright
  \item \emph{Say it:} \emph{āk}
  \item \emph{Say it:} \emph{āk}, once more
  \item \emph{Say it:} \emph{ret mā̃y āk hai.} — There's swallow-wort in the sand.
  \item \emph{Recall:} say the Marwadi for to go down (of the sun), then the Marwadi for the earth, then say \emph{āk} again
\end{itemize}

\begin{tcolorbox}[breakable,colback=teal!4,colframe=teal!35!black,title={Before you move on}]
What does \mw{आक} mean? (\textquotedblleft{}Swallow-wort (a desert shrub)\textquotedblright{}.) Say it once more.
\end{tcolorbox}
\section[\texorpdfstring{\mw{कांटो} (kā̃ṭo)}{कांटो (kā̃ṭo)}]{\mw{कांटो} (kā̃ṭo) — a thorn}

\label{lesson:MW-C230-kato}



\begin{quote}
Before the new one: say the Marwadi for the earth, then the Marwadi for swallow-wort.
\end{quote}

\subsection*{You'll want to know: \mw{कांटो}}
\textbf{\mw{कांटो}} — \emph{kā̃ṭo} — \textquotedblleft{}a thorn\textquotedblright{}.

\begin{quote}
\mw{पग} \mw{मांय} \mw{कांटो} \mw{है।}
\end{quote}

\begin{quote}
\emph{pag mā̃y kā̃ṭo hai.}
\end{quote}

\begin{quote}
There's a thorn in my foot.
\end{quote}

A masculine noun.

\begin{cousinweb}
\textbf{\mw{कांटो}} (\emph{kā̃ṭo}) is inherited from Sanskrit \emph{kaṇṭaka}, \textquotedblleft{}a thorn\textquotedblright{}; CDIAL lists Hindi \emph{kā̃ṭā} and a Gujarati cousin, \textquotedblleft{}a thorn\textquotedblright{}. CDIAL gives no Marwari form in this entry (CDIAL 2668).
\end{cousinweb}

\begin{grammarlens}[title={What you've built}]
One more word for this chapter.
\end{grammarlens}

\subsection*{Your turn}
\begin{itemize}
\raggedright
  \item \emph{Say it:} \emph{kā̃ṭo}
  \item \emph{Say it:} \emph{kā̃ṭo}, once more
  \item \emph{Say it:} \emph{pag mā̃y kā̃ṭo hai.} — There's a thorn in my foot.
  \item \emph{Recall:} say the Marwadi for the earth, then the Marwadi for swallow-wort, then say \emph{kā̃ṭo} again
\end{itemize}

\begin{tcolorbox}[breakable,colback=teal!4,colframe=teal!35!black,title={Before you move on}]
What does \mw{कांटो} mean? (\textquotedblleft{}A thorn\textquotedblright{}.) Say it once more.
\end{tcolorbox}
